Data centres allocate processors, memory and network bandwidth to thousands of jobs whose
demands are uncertain and change over time. Hand-designed scheduling heuristics work well for
the conditions their designers anticipated and poorly for the rest. This thesis studies
reinforcement learning as a way to learn allocation policies from experience. It formulates
cluster scheduling as a Markov decision process whose state encodes queued jobs and free
resources, trains policies with a graph-structured network that generalizes across cluster
sizes, and adds a constraint layer that guarantees no allocation exceeds capacity. In
simulations driven by production traces, the learned scheduler reduces average job completion
time by 19\% relative to the best heuristic while never violating a capacity constraint.
